\subsection{Secondary Robustness and Decidability Boundary}

The main side-information theorems address whatever collision structure is present in the realized representation. A complementary question is whether that fiber geometry is stable under system evolution: can a representation that is currently collision-free become noninjective when the class universe grows? The next two results are secondary robustness statements quantifying that fragility.

\leanmetapending{\LH{ROB1}}
\begin{definition}[Finite realized world]
Let $\mathcal{C}$ be the class universe. A \emph{finite realized world} is a finite subset $W \subseteq \mathcal{C}$ of classes currently present in a deployed system snapshot. We call $W$ \emph{barrier-free} iff
\[
\forall c_1,c_2 \in W,\ \pi_{\mathcal C}(c_1)=\pi_{\mathcal C}(c_2) \Rightarrow c_1=c_2.
\]
\end{definition}

\leanmetapending{\LHrng{ROB}{1}{2}}
\begin{definition}[Open-world extension]
An \emph{open-world extension} is a super-world $W' \supseteq W$ obtained by adding classes while keeping the observation family $\Phi$ fixed.
\end{definition}

\leanmetapending{\LHrng{ROB}{1}{2}}
\begin{theorem}[Barrier-freedom is not extension-stable]\label{thm:open-world-extension-instability}
In the open-world class model used here, it is not true that barrier-freedom is preserved under all extensions:
\[
\neg\Big(\forall W \subseteq W',\ \mathrm{BarrierFree}(W)\Rightarrow \mathrm{BarrierFree}(W')\Big).
\]
\end{theorem}

\begin{proof}
Take the empty world $W_0=\varnothing$, which is barrier-free. In this model, two concrete classes are shape-equivalent (same observable profile) but nominally distinct, so adding them yields an extension $W_1\supseteq W_0$ that is not barrier-free. Therefore barrier-freedom is not extension-stable.
\end{proof}

\begin{remark}[Safety interpretation]
Present collision-freedom does not guarantee stability under extension; persistent zero-error claims need an explicit growth model rather than a one-time snapshot argument.
\end{remark}

For compression system design: a representation needing no auxiliary bits today can require positive budget after extension. Open-world growth changes the optimal rate even when the deployed decoder is unchanged.

\leanmetapending{\LHrng{UND}{1}{2}}
\begin{definition}[Function-level barrier predicate]
For a partial observer generator $f:\mathbb{N}\rightharpoonup\mathbb{N}$, define
\[
\mathrm{HasBarrier}(f)\;:\Leftrightarrow\;\exists x\neq y,\exists z,\ z\in f(x)\land z\in f(y).
\]
Define $\mathrm{BarrierFree}(f):\Leftrightarrow \neg\mathrm{HasBarrier}(f)$.
\end{definition}

\leanmetapending{\LHrng{UND}{1}{2}}
\begin{theorem}[Rice-style non-computability of barrier certification]\label{thm:rice-barrier}
There is no computable predicate on program codes deciding whether the generated observer has a barrier:
\[
\neg\mathrm{ComputablePred}\!\left(c\mapsto \mathrm{HasBarrier}(\mathrm{eval}(c))\right).
\]
Equivalently, barrier-freedom certification is also non-computable.
\end{theorem}

\begin{proof}
The property is non-trivial (some generators produce collisions, some do not), so by Rice's theorem no computable predicate can decide it from code.
\end{proof}

\begin{remark}[Certification interpretation]
No general computable procedure can certify barrier-freedom for arbitrary future-generating code, even when deployment appears collision-free. Zero-error guarantees from finite training cannot extend to new production classes without explicit identity metadata or restrictive growth models.
\end{remark}

\leanmetapending{\LH{ROB1}, \LH{ROB2}, \LH{UND1}, \LH{UND2}}
\begin{corollary}[Certification consequence for persistent \texorpdfstring{$D=0$}{D=0} claims]
For unrestricted open-world generators, one cannot computably certify that attribute-only identification remains barrier-free under all future extensions. Barrier-freedom fails under arbitrary extensions \leanmeta{\LHrng{ROB}{1}{2}}, and both barrier-existence and barrier-freedom certification are uncomputable by Rice's theorem \leanmeta{\LHrng{UND}{1}{2}}. Consequently, persistent $D=0$ guarantees cannot rely on ``currently collision-free'' attribute structure alone.
\end{corollary}
